\frontmatterentry{To the reader}

\begin{enumerate}
\def\labelenumi{\arabic{enumi}.}
\item
  The Elements of Mathematics Series takes up mathematics at the beginning, and gives complete proofs. In principle, it requires no particular knowledge of mathematics on the readers' part, but only a certain familiarity with mathematical reasoning and a certain capacity for abstract thought. Nevertheless, it is directed especially to those who have a good knowledge of at least the content of the first year or two of a university mathematics course.
\item
  The method of exposition we have chosen is axiomatic, and normally proceeds from the general to the particular. The demands of proof impose a rigorously fixed order on the subject matter. It follows that the utility of certain considerations will not be immediately apparent to the reader unless he has already a fairly extended knowledge of mathematics.
\item
  The series is divided into Books and each Book into chapters. The Books already published, either in whole or in part, in the French edition, are listed below. When an English translation is available, the corresponding English title is mentioned between parentheses. Throughout the volume a reference indicates the English edition, when available, and the French edition otherwise.
\end{enumerate}

Théorie des Ensembles (Theory of Sets) designated by E (S)

Algèbre (Algebra \(^{(1)}\) ) --- A (A)

Topologie Générale (General Topology) --- TG (GT)

Fonctions d'une Variable Réelle --- FVR

Espaces Vectoriels Topologiques (Topological Vector Spaces) --- EVT (TVS)

Intégration --- INT

Algèbre Commutative (Commutative Algebra \(^{(2)}\) ) --- AC (CA)

Variétés Différentielles et Analytiques --- VAR

Groupes et Algèbres de Lie (Lie Groups and Lie Algebras \(^{(3)}\) ) --- LIE (LIE) Théories Spectrales --- TS

\pandocbounded{\includegraphics[keepaspectratio]{images/part001_af8bcccf.jpg}}

In the first six books (according to the above order), every statement in the text assumes as known only those results which have already been discussed in the same chapter, or in the previous chapters ordered as follows: S; A, chapters I to III; GT, chapters I to III; A, from chapters IV on; GT, from chapter IV on; FVR; TVS; INT.

From the seventh Book on, the reader will usually find a precise indication of its logical relationship to the other Books (the first six Books being always assumed to be known).

\begin{enumerate}
\def\labelenumi{\arabic{enumi}.}
\setcounter{enumi}{3}
\item
  However we have sometimes inserted examples in the text which refer to facts the reader may already know but which have not yet been discussed in the series. Such examples are placed between two asterisks: *\ldots*. Most readers will undoubtedly find that these examples will help them to understand the text. In other cases, the passages between *\ldots* refer to results which are discussed elsewhere in the Series. We hope the reader will be able to verify the absence of any vicious circle.
\item
  The logical framework of each chapter consists of the definitions, the axioms, and the theorems of the chapter. These are the parts that have mainly to be borne in mind for subsequent use. Less important results and those which can easily be deduced from the theorems are labelled as « propositions », « lemmas », « corollaries », « remarks », etc. Those which may be omitted at a first reading are printed in small type. A commentary on a particularly important theorem appears occasionally under the name of « scholium ».
\end{enumerate}

To avoid tedious repetitions it is sometimes convenient to introduce notations or abbreviations which are in force only within a certain chapter or a certain section of a chapter (for example, in a chapter which is concerned only with commutative rings, the word « ring » would always signify « commutative ring »). Such conventions are always explicitly mentioned, generally at the beginning of the chapter in which they occur.

\begin{enumerate}
\def\labelenumi{\arabic{enumi}.}
\setcounter{enumi}{5}
\item
  Some passages in the text are designed to forewarn the reader against serious errors. These passages are signposted in the margin with the sign \(\sum\) (« dangerous bend »).
\item
  The Exercises are designed both to enable the reader to satisfy himself that he has digested the text and to bring to his notice results which have no place in the text but which are nonetheless of interest. The most difficult exercises bear the sign \(\P\) .
\item
  In general, we have adhered to the commonly accepted terminology, except where there appeared to be good reasons for deviating from it.
\item
  We have made a particular effort always to use rigorously correct language, without sacrificing simplicity. As far as possible we have drawn attention in the text to abuses of language, without which any mathematical text runs the risk of pedantry, not to say unreadability.
\item
  Since in principle the text consists of the dogmatic exposition of a theory, it contains in general no references to the literature. Bibliographical references are gathered together in Historical Notes. The bibliography which follows each historical note contains in general only those books and original memoirs which have been of the greatest importance in the evolution of the theory under discussion. It makes no sort of pretence to completeness.
\end{enumerate}

As to the exercises, we have not thought it worthwhile in general to indicate their origins, since they have been taken from many different sources (original papers, textbooks, collections of exercises).

\begin{enumerate}
\def\labelenumi{\arabic{enumi}.}
\setcounter{enumi}{10}
\tightlist
\item
  In the present Book, references to theorems, axioms, definitions, \ldots{} are given by quoting successively :
\end{enumerate}

--- the Book (using the abbreviation listed in Section 3), chapter and page, where they can be found, when referring to the French edition;

--- the chapter and page only when referring to the present Book;

--- the chapter, paragraph and section, when referring to the English edition.

The Summaries of Results are quoted by the letter R; thus Set Theory, R signifies « Summary of Results of the Theory of Sets ».

